\documentclass[course]{neyvia-study}
\studytitle{Linear Algebra}
\studysubtitle{Complete Course}
\studytagline{Example handout \ ·\ one chapter}
\studysourceprefix{handout }
\begin{document}
\maketitle

\studychapter{Sets and maps}
\studypart{I}{Sets}

\begin{concept}{Set equality}{p1 col1}
\begin{definition}
Let $E$ be a set and $A,B\subseteq E$.
\dline{A=B\iff\big(\forall x\in E,\ x\in A\iff x\in B\big)}
\end{definition}
\begin{method}{Prove that $A=B$ (double inclusion)}
\step{\kw{$\subseteq$)} \kw{Let} $x\in A$. $\ \cdots\ $ so $x\in B$.}
\step{\kw{$\supseteq$)} \kw{Let} $x\in B$. $\ \cdots\ $ so $x\in A$.}
\step{$\therefore A=B$.}
\end{method}
\begin{reflex}
Set equality \ra{} prove both inclusions.
\end{reflex}
\begin{worked}{Two descriptions of one set}
\step{Let $A=\set{x\in\R\mid x^2=1}$ and $B=\set{-1;1}$.}
\step{$\displaystyle x\in A\iff x^2=1$}
\step{$\displaystyle x\in A\iff x=-1\ \text{or}\ x=1$}
\result{A=B}
\end{worked}
\begin{traps}
\trap{$\set{1;2}=\set{2;1}$ but $(1;2)\neq(2;1)$.}
\end{traps}
\begin{practice}
\try{Is $\set{x\in\Z\mid x^2<1}=\set{0}$?}
\answer{Yes: the only integer with $x^2<1$ is $0$.}
\end{practice}
\studynote{The order of elements in a set never matters.}
\end{concept}

\begin{concept}{Surjective maps}{p1 col2}
\begin{definition}
Let $f:E\to F$ be a map.
\dline{f \text{ surjective} \iff \forall y\in F,\ \exists x\in E,\ y=f(x)}
\end{definition}
\begin{method}{Prove that $f$ is surjective}
\step{Let $y\in F$.}
\step{We look for $x\in E$ such that $f(x)=y$.}
\step{Take $x=\cdots$ (it lies in $E$).}
\step{Then $f(x)=\cdots=y$.}
\step{$\therefore f$ is surjective.}
\end{method}
\begin{reflex}
Surjective \ra{} start with $y\in F$, solve $y=f(x)$.
\end{reflex}
\end{concept}
\end{document}
